\chapter*{Anexo D: Repositorio del proyecto y guía de uso}
\addcontentsline{toc}{chapter}{Anexo D: Repositorio del proyecto y guía de uso}
\label{cap:anexo_repositorio}

El código fuente completo de Open CVN, junto con esta misma memoria, está
disponible públicamente en el siguiente repositorio:

\begin{center}
\url{https://github.com/holt00/TFG-open-cvn-schema}
\end{center}

Este anexo describe la organización general del repositorio y ofrece una
guía de orientación para quien quiera consultarlo o ejecutarlo, sin repetir
el detalle de comandos ya recogido en el Anexo A.

\section*{Estructura del repositorio}

El repositorio se organiza en un pequeño número de directorios de nivel
superior, cada uno con una responsabilidad clara:

\begin{itemize}
  \item \path{src/}: código fuente del proyecto, organizado según la
  arquitectura por capas descrita en el Capítulo~\ref{cap:metodologia}:
  bindings estructurales generados, lógica de generación y normalización
  mantenida a mano, modelos de dominio, el contrato público de análisis y
  validación, y la herramienta local de gestión.

  \item \path{schemas/}: artefacto JSON Schema canónico del formato Open CVN,
  detallado en el Anexo B.

  \item \path{examples/}: documentos Open CVN JSON de ejemplo, citados a lo
  largo de esta memoria.

  \item \path{tests/}: batería de pruebas automatizadas evaluada en el
  Capítulo~\ref{cap:evaluacion}.

  \item \path{docs/}: documentación técnica del proyecto, organizada por
  tema (arquitectura, pipeline, formato, desarrollo) y el registro de
  planificación por issues y correcciones puntuales resumido en el Anexo C.
  El propio código fuente de esta memoria se encuentra en
  \path{docs/memoria/}.
\end{itemize}

\section*{Cómo orientarse en el repositorio}

Para quien no haya seguido el desarrollo del proyecto, el punto de entrada
recomendado es el fichero \path{README.md} en la raíz del repositorio, que
resume el objetivo del proyecto y su alcance actual. El fichero
\path{PROJECT_GUIDE.md} amplía esa descripción con un mapa completo de la
documentación técnica disponible bajo \path{docs/}, incluida la
correspondencia entre cada capítulo de esta memoria y su documentación
técnica de referencia en el repositorio.

Para preparar el entorno, regenerar los artefactos y ejecutar la batería de
pruebas o la herramienta local, el Anexo A recoge los comandos exactos y un
recorrido completo de ejemplo.
